% DRAFT generated by akm.run draft. Keep what applies; add anything triage revealed.
\begin{itemize}
  \item \emph{Crash model.} F6 kills the migrator process; it does not cut power. Data the operating system
        had buffered but not yet flushed survives a process crash but might not survive power loss.
  \item \emph{Platform.} Timings come from one machine. Under WSL2, fsync passes through a virtual disk, so
        absolute times and durability are platform-specific; ratios are more portable than seconds.
  \item \emph{Warm caches.} Caches were not dropped between repetitions; a discarded warm-up run keeps the
        measured repetitions under the same conditions.
  \item \emph{Isolation strength.} The file-access guard used for bulk runs prevents accidental access by our
        own code, not a deliberately malicious verifier; the Docker configuration is the stronger layer.
  \item \emph{Shared codebase.} The oracle V3 shares code with the verifiers it grades; every trial therefore
        cross-checks V3 against the independent fault manifest, and any disagreement halts the run.
  \item \emph{Synthetic data and seeded keys.} Records are synthetic and keys are derived from a seed rather
        than the operating system's random source. Neither affects detection; content could slightly
        affect timing.
  \item \emph{Sampling analysis.} Each RQ4 challenge sequence serves every $k$ as a prefix, so points on one
        curve are correlated; the pass/fail decision uses a Bonferroni correction, which remains valid
        under such dependence.
  \item \emph{Single group-commit size.} Group commit was measured at $B=1000$ only.
  \item \emph{Excluded large-scale configurations.} Experiment~2 excluded $N=100{,}000$ and the 32{,}768\,B
        payload size from per-record-fsync timing, since measured per-write fsync cost at that record size
        (6.98\,ms, versus 3.15\,ms for small writes) made the full grid impractical under WSL2; storage
        measurements at these sizes were still collected where they do not require per-record timing.
  \item \emph{Confidence-interval fragility near saturation.} Wilson intervals narrow sharply near 0\% or
        100\% detection; the one RQ4 cell falling outside its Bonferroni-adjusted interval (F3,
        $\varepsilon=5\%$, $k=200$) did so at a 99.9\% observed rate against a 99.997\% prediction, a gap
        attributable to this known statistical property rather than a detection failure.
\end{itemize}
